\documentclass{beamer}
\input{semantics.tex}
\title[Combinations of Words]{Words in Combination: Idioms and Composition}

\usepackage{forest}
\usepackage{tikz}
\usepackage{tikz-dependency}
\newcommand{\sr}[1]{\ensuremath{\langle}#1\ensuremath{\rangle}}


\begin{document}
 
\maketitle
\input{part/word-meaning-revisited} 
\input{part/composition}
\section{Multi-Word Expressions}
\newcommand{\tra}[1]{\textcolor{olive}{\textsf{#1}}}



\begin{frame}{Idioms}

\begin{itemize}
\item Some expressions clearly involve more than one orthographic word
  \begin{itemize}
  \item compound noun
    \begin{itemize}
    \item \eng{grass snake}; \eng{grass and tree snakes}
    \end{itemize}
  \item verb-particle
    \begin{itemize}
    \item \eng{I looked it up} vs \eng{I looked up the very long word}
    \end{itemize}
  \item  idiom
    \begin{itemize}
    \item \eng{going great guns, give the Devil his due}
    \item \eng{jog someone's memory}
    \item \eng{blow one's top}
    \end{itemize}
\item And more
\\ \eng{San Francisco, ad hoc, by and large, Where Eagles
Dare, kick the bucket, part of speech, in step, the
Oakland Raiders, trip the light fantastic, telephone
box,  take a walk , do a number on
(someone), take (unfair) advantage (of), pull strings,
kindle excitement, fresh air, \ldots}
  \end{itemize}
\item Knowing the individual words is not enough to know the meaning (or usage)
  
\end{itemize}

\end{frame}

\begin{frame}{Multiword Expressions (MWE)}

  There are many different kinds of irregularity.\\[2ex]
%  \bigskip
\begin{small}
\begin{tabular}{lccccc}
  MWE & \multicolumn{5}{c}{Weirdness} \\
                  & Lexical  & Syntactic & Semantic & Pragmatic &
                                                                  Statistical  \\
\hline
\eng{ad hominem}       &  $+$ & ?   &  ?   &   ?   & $+$ \\
\eng{at first}          &      & $+$ &     &      & ? \\
\eng{first aid}         &      &      & $+$   &      & ? \\
\eng{salt and pepper}   &      &      &      &      & $+$\\
\eng{good morning}      &      &      &      &  $+$ &$+$ \\
\eng{cat's cradle}     & $+$  &      &  $+$ &      & ?\\
\end{tabular}
  \end{small}
\begin{itemize}
\item Most of the time, we don't even notice
\item Unless it is your second language
% \item In Project Three you will find examples of interesting
%   multi-word expressions
% \\ from a new story
\item Much more in \citet{Sag02a}
\end{itemize}

\end{frame}

\begin{frame}{We analyze these differently}

  \begin{itemize}
  \item \txx{Fixed Expressions} are like words with spaces
    \\ \lex{ad hominem} ``at the person''
  \item \txx{Semi-Fixed Expressions} allow some inflection
    \\ \lex{part(s)-of-speech} ``syntactic category''
  \item \txx{Flexible Expressions} \citep{Nunberg:Sag:Wasow:1994}
    \begin{itemize}
    \item \txx{Decomposable Expressions} we can treat the parts as
      having different meanings in context
      \\ \lex{spill the beans} ``reveal the secrets''
    \item \txx{Non-decomposable Expressions}
      \\ \lex{kick the bucket} ``die''
    \end{itemize}
  \item \txx{Institutionalized Phrases} are statistically marked
    \\ \lex{black and white} vs $^?$\lex{white and black} c.f. \lex{白黒} \jpn[white black]{shirokuro}
    \\ \lex{machine translation} vs $^?$\lex{computer translation}
    
  \end{itemize}
  
\end{frame}

\begin{frame}{Semantic Opacity: A Gradient}

\begin{exe}
  \ex \textbf{Opaque:} \emph{kick the bucket} $\Rightarrow$ meaning unrelated to kicking/buckets
  \ex \textbf{Partially transparent:} \emph{spill the beans} $\Rightarrow$ metaphorical extension
  \ex \textbf{Highly transparent:} \emph{strong tea} $\Rightarrow$ collocational preference
\end{exe}


\vspace{0.5em}
Dimensions of semantic idiosyncrasy:
\begin{itemize}
  \item Non-compositional meaning
  \item Restricted argument structure
  \item Fixed lexical choices
  \item Conventional metaphor/metonymy
\end{itemize}
\end{frame}

% ------------------------------------------------------------------------
\begin{frame}{Verb-Particle Constructions}

\begin{exe}
  \ex \eng{She looked up the word}
  \ex \eng{She looked the word up}
\end{exe}

\vspace{0.5em}
Particles may:
\begin{itemize}
  \item Encode direction (\eng{walk out})
  \item Encode aspect (\eng{eat up})
  \item Encode idiomatic meaning (\eng{give up})
\end{itemize}
\end{frame}

% ------------------------------------------------------------------------
\begin{frame}{Collocations and Selectional Preferences}
\begin{exe}
  \ex \eng{strong tea} vs.\ *\eng{powerful tea}
  \ex \eng{heavy rain} vs.\ *\eng{strong rain}
\end{exe}

\vspace{0.5em}
Not fully idiomatic, but statistically entrenched.

\vspace{0.5em}
For computational semantics:
\begin{itemize}
  \item Collocations affect vector representations.
  \item Lexical choice is conventionalized.
\end{itemize}
\end{frame}

% ------------------------------------------------------------------------
\begin{frame}{Formulaic Sequences and Lexical Bundles}

  In corpus linguistics \citep{Biber:Conrad:Reppen:1998}:

\begin{exe}
  \ex \eng{on the other hand}
  \ex \eng{as a result of}
  \ex \eng{in the case of}
\end{exe}

\vspace{0.5em}
Properties:
\begin{itemize}
  \item High frequency.
  \item Often discourse-organizing.
  \item Not necessarily semantically opaque.
\end{itemize}

\vspace{0.5em}
Bundles show how usage frequency shapes grammar.
\end{frame}

\begin{frame}{How common are MWEs?}

\begin{itemize}
\item They are very common in the lexicon
  \begin{itemize}
  \item In wordnet,  41\% of the entries are multiword (mainly compound nouns)
  \end{itemize}
\item But less common in the actual text (\sh{SPEC} 4.5\%: 296/6,641)
\\ 24 are new (not in Wordnet 3.0); 55 are named entities
% sqlite3 /var/www/ntumc/db/eng.db "select substr(tag,1,1), clemma, tag, count(clemma) from concept where sid >=11000 and sid <=11700 and tag != 'x' and clemma glob '* *' group by substr(tag,1,1)"
% -- Loading resources from /home/bond/.sqliterc

% substr(tag,1,1)	clemma	tag	count(clemma)
% 0	and so	00117620-r	206
% 1	good fortune	11463746-n	11

% 7	no one	77000137-n	3
% 8	pull off	80000623-v	16
% 9	last night	90000518-n	5

% d	four o'clock	dat	1
% l	baker street	loc	18
% o	in order	oth	3
% p	Hilton Cubitt	per	32
% w	as well as	w	1

  \begin{itemize}
  \item \eng{take into one's confidence}
  \item \eng{take in}
  \item \eng{Sherlock Holmes}
  \item \eng{practical joke(r)}
  \item \eng{in love}
  \item \eng{get the better of}
  \item \eng{Panama hat}
  \item \eng{as good as one's word} 
  \end{itemize}
\item It still seems as though wordnet (and dictionaries in general) are missing many MWEs
\end{itemize}

\end{frame}

\begin{frame}{Why are they important?}

\begin{itemize}
\item If you think you know the individual words, then you might be  confused

\item  Which is a problem if you are a translator:
\\ \eng[whoever he met]{whoever crossed his path} \sh{SPEC}
\begin{exe}
  \ex \glll 私道 を 渡ろう と する 人 \\
  shidou wo watarou to suru hito \\
  private-road \textsc{acc} cross.let's \textsc{quot} do person \\
\trans ''whoever tried to cross his private road''
\end{exe}
\item Knowledge of MWEs is one of the things that separates a good speaker from a poor one
\item From a linguist's point of view, they also reveal something about
  how language is organized in our brains
\end{itemize}



%\section{Idioms}

%\begin{frame}{Idioms are hard to handle}
\end{frame}

\begin{frame}{Idioms are hard to handle}
%\MyLogo{They don't translate well}
%\MyLogo{\eng{atama-ni kuru} better with 1st person subject}
\begin{exe}
  \ex \eng{Kim blew her top}
  \trans ``Kim got angry''
  \trans No \emp{blowing}, no \emp{top}, no \emp{her}
  \ex  \glll \jpn{キムは} \jpn{頭に} \jpn{来た}  \\
  kimu-wa atama-ni  kita \\
  Kim-\textsc{top} head-\textsc{dat} came \\
  \trans ``Kim got angry (lit: Kim came to her head)''
  \trans no \emp{head}, no \emp{coming}
\end{exe}

\begin{itemize}
\item They are hard to identify
\item They require work to represent
\end{itemize}
%\end{frame}

%\begin{frame}{The state of the art (translation)}
\end{frame}

\begin{frame}{The state of the art (translation)}
\MyLogo{Not so good}

\begin{exe}
  \ex \eng{Kim blew her top}
  \trans \glll \tra{キムは、} \tra{彼女の} \tra{上を} \tra{吹きました} \\
  {Kimu wa,} {kanojo no} {ue o} fukimashita \\
  Kim-\textsc{top} her-'s top-\textsc{acc} blew \\  
  \trans  \tra{Kim exhaled on the upper part of her} (Google Translate
  2015)
 \trans \glll \tra{キムは、} \tra{彼女の} \tra{トップを} \tra{吹いた} \\
  {Kimu wa,} {kanojo no} {toppu o} fuita \\
  Kim-\textsc{top} her-'s top-\textsc{acc} blew \\  
  \trans  \tra{Kim exhaled on the shirt of her} (Bing Translate 2015)
 \trans \glll \tra{キムは、} \tra{トップを} \tra{吹いた} \\
  {Kimu wa,} {toppu o} fuita \\
  Kim-\textsc{top}  top-\textsc{acc} blew \\  
  \trans  \tra{Kim exhaled on the shirt} (Google Translate 2019)
  \ex \eng{キムは 頭に 来た}
  \trans \tra{Kim came to the head} (Google+Bing Translate 2015)
%  \trans \tra{Kim came to the head} (Bing Translate)
\end{exe}

\end{frame}

\begin{frame}{The state of the art (wordnet)}
\MyLogo{Some but not all}
\vspace*{-1.5ex}
\begin{itemize}
\item Hyponym of \eng{rage} with many lemmas
\begin{multicols}{2}
  \begin{itemize}
  \item \eng{flip one's lid}
  \item \eng{blow one's stack}
  \item \eng{flip one's wig}
  \item \eng{lose one's temper}
  \item \eng{blow up}
  \item \eng{throw a fit}
  \item \eng{hit the roof}
  \item \eng{hit the ceiling}
  \item \eng{have kittens}
  \item \eng{have a fit}
  \item \eng{combust}
  \item \eng{fly off the handle}
  \item \eng{blow a fuse}
  \item \eng{go ballistic}
  \end{itemize}
\end{multicols}
\item But by no means exhaustive
\begin{multicols}{2}
  \begin{itemize}
  \item \eng{lose one's cool}
  \item \eng{lose one's rag}  
  \item \eng{lose one's grip}
  \item \eng{blow one's top}
  \item \eng{fly off one's handle}
  \item \eng{spit the dummy}
  \end{itemize}
\end{multicols}
\end{itemize}

% \end{frame}

\end{frame}

\begin{frame}{The state of the art (large language models)}
% \MyLogo{Actually quite good}
% %\makexeCJKactive
  \begin{small}
    \begin{itemize}
    \item FCB: Translate "Kim blew her top" into Japanese
    \item ChatGPT:
      \texttt{The English idiom "Kim blew her top" translates to Japanese as ``キムはキレた" (Kimu wa kireta). This phrase captures the sense that Kim lost her temper or became very angry.}
    \item FCB: What does ``キムは頭に来た" mean
    \item ChatGPT:  \texttt{The phrase ``キムは頭に来た" (Kimu wa atama ni kita) translates to English as "Kim got angry" or more literally, "It came to Kim's head." It's another way to express that Kim became angry or upset, similar to saying "Kim blew her top" in English. \\
        The expression  ``頭に 来た" (atama ni kita) is a common way in Japanese to describe someone getting angry or frustrated.}
    \end{itemize}
  \end{small}
  % \makexeCJKinactive
\hfill Translated with openai's ChatGPT 4.0 (2024-04-21)

Training on a lot of data, with a lot of parameters, makes it much better

\end{frame}


\input{part/possessed-idioms}
%\section{Constructions}

\begin{frame}{Beyond words and phrases}

\begin{itemize}
\item Some uses of language seem to come from beyond just the words
  \begin{itemize}
  \item \eng{They laughed the poor guy out of the room.}
  \item \eng{Mary urged Bill into the house.}
  \end{itemize}
\item \textbf{Caused Motion Construction}
  \\ CAUSE-MOVE $<$agent, theme, goal/path$>$ \into S V NP PP
  \begin{itemize}
  \item compatible verbs can fuse with the construction
  \end{itemize}
\item \textbf{The Xer the Yer}
  \begin{itemize}
  \item \eng{the more you think about it, the less you understand}
  \end{itemize}
\item At one end, phrases can be modeled as very general constructions
\item At the other end, idioms can be modeled as very specific constructions
\end{itemize}

\end{frame}

% --- Forest style for construction networks -----------------------------
\forestset{
  cxg/.style={
    for tree={
      grow'=east,
      parent anchor=east,
      child anchor=west,
      align=left,
      rounded corners,
      draw,
      inner sep=2pt,
      s sep=6pt,
      l sep=10pt,
      edge={-latex}
    }
  }
}

% ------------------------------------------------------------------------
\begin{frame}{Roadmap (20 slides)}
\begin{itemize}
  \item Start with \textbf{non-canonical constructions}: \emph{whistle your way to work}
  \item Key idea: \textbf{form--meaning pairings} at multiple granularities
  \item From ``weird'' to \textbf{argument structure}, \textbf{idioms}, \textbf{morphology}, \textbf{discourse}
  \item Semantics focus: \textbf{coercion}, \textbf{frames}, \textbf{constructional meaning}
  \item How to represent: \textbf{AVMs}, \textbf{inheritance networks}, \textbf{examples}
\end{itemize}
\end{frame}

% ------------------------------------------------------------------------
\begin{frame}{Warm-up: a non-canonical pattern}
\small
% \label{ex:way1}
\begin{exe}
  \ex \gll She \textbf{whistled} her way to work.\\
      she whistle.\textsc{pst} her way to work\\
      \glt `She went to work while whistling (and the whistling is part of how she progressed).'
  \ex \gll He \textbf{elbowed} his way through the crowd.\\
      he elbow.\textsc{pst} his way through the crowd\\
      \glt `He made progress through the crowd by using his elbows.'
  \ex \gll They \textbf{googled} their way to an answer.\\
      they google.\textsc{pst} their way to an answer\\
      \glt `They reached an answer by googling (repeatedly / strategically).'
\end{exe}


\vspace{0.5em}
\textbf{Observation:} the verb contributes \emph{manner/means}, but the construction contributes \emph{directed motion + incremental progress}.
\end{frame}

% ------------------------------------------------------------------------
\begin{frame}{Why this is interesting (semantics)}
\small
\begin{itemize}
  \item The verb \emph{whistle} is not a motion verb.
  \item Yet the sentence expresses a \textbf{motion event} with a \textbf{path/goal}.
  \item Meaning is not (just) verb-centered; it is \textbf{construction-centered}.
\end{itemize}

\vspace{0.75em}
\textbf{Constructional meaning sketch:}
\begin{itemize}
  \item \textsc{event}: motion-to-goal
  \item \textsc{means}: activity denoted by V
  \item \textsc{incrementality}: progress measured by ``way'' path
\end{itemize}
\end{frame}

% ------------------------------------------------------------------------
\begin{frame}{CxG in one sentence}
\large
\begin{center}
\textbf{A construction is a conventionalized pairing of form and meaning}\\
\textbf{(including constraints on usage).}
\end{center}

\vspace{0.75em}
\normalsize
\begin{itemize}
  \item Constructions range from \textbf{morphemes} to \textbf{idioms} to \textbf{abstract patterns}.
  \item Grammar is a \textbf{network} (inheritance, similarity, family resemblance).
\end{itemize}
\end{frame}

% ------------------------------------------------------------------------
\begin{frame}{From “words + rules” to a constructicon}
\small
Traditional modular view (oversimplified):
\begin{itemize}
  \item Lexicon: words (meanings)
  \item Syntax: rules (structures)
\end{itemize}

\vspace{0.5em}
CxG view:
\begin{itemize}
  \item Lexicon and syntax are a \textbf{continuum} of constructions.
  \item The inventory of constructions = \textbf{constructicon}.
\end{itemize}

\vspace{0.75em}
\textbf{Continuum examples:}

\begin{exe}
  \ex \emph{-ness} (morpheme construction)
  \ex \emph{kicked the bucket} (idiom)
  \ex \emph{the X-er, the Y-er} (comparative correlative)
  \ex \emph{NP V NP PP} (argument structure pattern)
\end{exe}

\end{frame}

% ------------------------------------------------------------------------
\begin{frame}{A first AVM: the Way Construction (simplified)}
% \small
% \avm{
% [ \textsc{cxn} & \textit{Way-Construction} \\
%   \textsc{form} & [ \textsc{syn} & [ \textsc{cat} & \textit{VP} \\
%                                   \textsc{head} & [ \textsc{v} & V ] \\
%                                   \textsc{comps} & \langle NP\textsubscript{poss}~\textit{way},~PP\textsubscript{path} \rangle ] ] \\
%   \textsc{meaning} & [ \textsc{event} & \textit{motion} \\
%                     \textsc{path} & \textit{directed} \\
%                     \textsc{means} & \textit{activity}(V) \\
%                     \textsc{progress} & + ] \\
%   \textsc{constraints} & [ \textsc{v-semantics} & \textit{manner/means} \\
%                         \textsc{pp} & \textit{path/to/through/into}\ldots ] ]
% }

\vspace{0.5em}
\textbf{Takeaway:} a \emph{template} carries meaning and selects compatible verbs/PPs.
\end{frame}

% ------------------------------------------------------------------------
\begin{frame}{Not just “weird”: Constructional coercion}
\small
 \label{ex:coercion}
\begin{exe}
  \ex \gll She \textbf{sneezed} the napkin off the table.\\
      she sneeze.\textsc{pst} the napkin off the table\\
      \glt `Her sneezing caused the napkin to move off the table.'
  \ex \gll He \textbf{laughed} the meeting into chaos.\\
      he laugh.\textsc{pst} the meeting into chaos\\
      \glt `His laughing caused the meeting to become chaotic.'
\end{exe}


\vspace{0.5em}
\textbf{Idea:} a construction can \textbf{coerce} an event into a causal / change-of-state / motion construal.
\end{frame}

% ------------------------------------------------------------------------
\begin{frame}{The Caused-Motion Construction}
\small
Canonical-ish schema:
 \label{ex:causedmotion}
\begin{exe}
  \ex \gll She pushed the box into the hallway.\\
      she push.\textsc{pst} the box into the hallway\\
      \glt `She caused the box to move into the hallway.'
  \ex \gll The trainer threw the dolphin through the hoop.\\
      the trainer throw.\textsc{pst} the dolphin through the hoop\\
      \glt `The trainer caused the dolphin to move through the hoop.'
  \ex \gll She sneezed the napkin off the table.\\
      she sneeze.\textsc{pst} the napkin off the table\\
      \glt `She caused the napkin to move off the table.'
\end{exe}


\vspace{0.5em}
\textbf{Same constructional meaning; different verbs.}
\end{frame}

% ------------------------------------------------------------------------
\begin{frame}{Argument structure as constructions}
\small
Instead of ``the verb projects its arguments'', CxG often treats patterns as meaningful.


\begin{exe}
  \ex \textbf{Ditransitive:} \emph{NP V NP NP} $\Rightarrow$ \textsc{transfer} / \textsc{caused possession}\\
      \gll She baked him a cake.\\
          she bake.\textsc{pst} him a cake\\
      \glt `She caused him to have a cake (by baking).'
  \ex \textbf{Resultative:} \emph{NP V NP AP/PP} $\Rightarrow$ \textsc{caused change-of-state}\\
      \gll They hammered the metal flat.\\
          they hammer.\textsc{pst} the metal flat\\
      \glt `They caused the metal to become flat (by hammering).'
  \ex \textbf{Intransitive motion:} \emph{NP V PP} $\Rightarrow$ \textsc{directed motion}\\
      \gll The bottle floated into the cave.\\
          the bottle float.\textsc{pst} into the cave\\
\end{exe}

\end{frame}

% ------------------------------------------------------------------------
\begin{frame}{Ditransitive: subtle meaning contrasts}
\small
Compare two ways to encode transfer/possession:


\begin{exe}
  \ex  She gave the student a book.\\
      \glt `Direct transfer; recipient is foregrounded.'
  \ex  She gave a book to the student.\\
      \glt `Path/goal phrasing; can sound more ``delivery-like''.'
\end{exe}


\vspace{0.75em}
\textbf{CxG angle:} these are \emph{distinct constructions} with overlapping but non-identical semantics and discourse profiles.
\end{frame}

% ------------------------------------------------------------------------
\begin{frame}{Construction networks (inheritance / family resemblance)}
\small
\begin{forest} cxg
[Argument-Structure Constructions
  [Motion-family
    [Intransitive Motion\\\emph{NP V PP}]
    [Caused Motion\\\emph{NP V NP PP}]
    [Way-Construction\\\emph{NP V NP\textsubscript{way} PP}]
  ]
  [Change-of-state family
    [Resultative\\\emph{NP V NP AP/PP}]
    [Into-Causative\\\emph{NP V NP into NP}]
  ]
  [Transfer family
    [Ditransitive\\\emph{NP V NP NP}]
    [To-Dative\\\emph{NP V NP to NP}]
  ]
]
\end{forest}

\vspace{0.5em}
\textbf{Key point:} generalizations live in the network, not only in ``rules''.
\end{frame}

% ------------------------------------------------------------------------
\begin{frame}{Idioms are constructions too (and they generalize)}
\small
Fixed idioms:

\begin{exe}
  \ex \emph{kick the bucket} `die'
  \ex \emph{spill the beans} `reveal a secret'
\end{exe}


\vspace{0.5em}
Partially schematic idioms (``open slots''):

\begin{exe}
  \ex \emph{the X-er, the Y-er}: \emph{The more you read, the happier you get.}
  \ex \emph{What’s X doing Y?}: \emph{What’s that dog doing in my office?}
  \ex \emph{let alone}: \emph{He can’t boil an egg, let alone cook dinner.}
\end{exe}


\vspace{0.5em}
\textbf{Semantics lives at the pattern level.}
\end{frame}

% ------------------------------------------------------------------------
\begin{frame}{Information structure constructions}
\small
Constructions can encode discourse-pragmatics.


\begin{exe}
  \ex \textbf{Cleft:} \emph{It was \textbf{John} who broke the vase.}\\
      Focuses \textbf{John} as the salient alternative.
  \ex \textbf{Topicalization:} \emph{That book, I haven’t read \_\ yet.}\\
      Promotes topic; leaves a gap.
  \ex \textbf{Existential:} \emph{There’s a spider in the bath.}\\
      Introduces a new discourse referent.
\end{exe}

\end{frame}

% ------------------------------------------------------------------------
\begin{frame}{Constructions in morphology and word formation}
\small
Not only sentences:


\begin{exe}
  \ex \textbf{-ness} construction: \emph{happy} $\to$ \emph{happiness}\\
      Meaning: abstract property/state
  \ex \textbf{X-\emph{gate}} construction: \emph{Watergate} $\to$ \emph{Emailgate}\\
      Meaning: scandal associated with X
  \ex \textbf{N--N compounds}: \emph{chicken soup}, \emph{stone wall}\\
      Meaning: a relation that must be inferred (ingredient / material / purpose / \ldots)
\end{exe}


\vspace{0.5em}
\textbf{CxG claim:} these too are stored, learned, and productive patterns.
\end{frame}

% ------------------------------------------------------------------------
\begin{frame}{Semantic frames: linking constructions to world knowledge}
\small
Many CxG approaches connect constructions to frame-like semantics.


\begin{exe}
  \ex \emph{Ditransitive} evokes a \textsc{transfer} / \textsc{caused-possession} frame.
  \ex \emph{Caused-motion} evokes a \textsc{cause-to-move} frame.
  \ex \emph{Way-construction} evokes a \textsc{directed-motion-with-means} frame.
\end{exe}


\vspace{0.75em}
\textbf{Why it matters for semantics:}
\begin{itemize}
  \item inference (what follows from the construal)
  \item roles (participant structure beyond syntax)
  \item alternations (meaning shifts with form)
\end{itemize}
\end{frame}

% ------------------------------------------------------------------------
\begin{frame}{From special cases to “everything is a construction”}
\small
The punchline:
\begin{itemize}
  \item If non-canonical patterns need stored form--meaning pairings,
  \item and partially productive patterns do too,
  \item and alternations correlate with meaning/discourse differences,
\end{itemize}
\vspace{0.5em}
\begin{center}
\Large
\textbf{then general grammar is best modeled as a network of constructions.}
\end{center}
\end{frame}

% ------------------------------------------------------------------------
\begin{frame}{How do constructions get learned? (usage-based intuition)}
\small
Usage-based story (common in CxG traditions):
\begin{itemize}
  \item Learners store exemplars; generalizations emerge from frequency and similarity.
  \item High-frequency patterns become entrenched.
  \item Productivity is graded (not all-or-nothing).
\end{itemize}

\vspace{0.75em}
Mini illustration:

\begin{exe}
  \ex Heard often: \emph{give NP NP}, \emph{send NP NP}, \emph{show NP NP} $\Rightarrow$ strong ditransitive schema
  \ex Heard rarely: \emph{whisper NP NP} $\Rightarrow$ marginal but interpretable via schema + pragmatics
\end{exe}

\end{frame}

% ------------------------------------------------------------------------
\begin{frame}{Cross-linguistic perspective (quick taste)}
\small
CxG is friendly to cross-linguistic comparison because you can compare \emph{constructions}, not only categories.


\begin{exe}
  \ex English: \emph{She laughed him out of the room.} (caused-motion coercion)
  \ex Some languages prefer: periphrasis / explicit causatives / serial verbs (language-specific constructions)
  \ex Topic-prominent languages often have rich topic constructions (form + discourse meaning)
\end{exe}


\vspace{0.5em}
\textbf{Semantics angle:} similar communicative functions may be packaged in different constructions.
\end{frame}

% ------------------------------------------------------------------------
\begin{frame}{Worked micro-analysis: Resultatives}
\small

\begin{exe}
  \ex \gll They wiped the table clean.\\
      they wipe.\textsc{pst} the table clean\\
      \glt `They caused the table to become clean (by wiping).'
  \ex \gll He drank the pub dry.\\
      he drink.\textsc{pst} the pub dry\\
      \glt `He caused the pub to become dry/empty (by drinking).'
\end{exe}


\vspace{0.5em}
\textbf{Notes:}
\begin{itemize}
  \item The adjective encodes a \textbf{result state}.
  \item Some pairs are natural; others require imaginative bridging/inference.
  \item Productivity is constrained by plausibility and convention.
\end{itemize}
\end{frame}

% ------------------------------------------------------------------------
\begin{frame}{Representing a construction as an AVM (Resultative, simplified)}
\small
% \avm{
% [ \textsc{cxn} & \textit{Resultative-Cxn} \\
%   \textsc{form} & [ \textsc{syn} & [ \textsc{cat} & \textit{VP} \\
%                                   \textsc{head} & V \\
%                                   \textsc{comps} & \langle NP\textsubscript{obj},~XP\textsubscript{res} \rangle ] ] \\
%   \textsc{meaning} & [ \textsc{event} & \textit{cause} \\
%                     \textsc{theme} & NP\textsubscript{obj} \\
%                     \textsc{result} & \textit{state}(XP\textsubscript{res}) ] \\
%   \textsc{constraints} & [ \textsc{compatibility} & \textit{causal plausibility} ] ]
% }

\vspace{0.5em}
\textbf{Use in semantics teaching:} make explicit what is contributed by pattern vs. lexeme.
\end{frame}

% ------------------------------------------------------------------------
\begin{frame}{Mini “constructional minimal pairs”}
\small
Try these as in-class discussion prompts.


\begin{exe}
  \ex \emph{She ran to the station.} (directed motion)
  \ex \emph{She ran the shoes to pieces.} (resultative change-of-state)
  \ex \emph{She ran her way to the station.} (way construction)
\end{exe}


\vspace{0.5em}
\textbf{Question:} what semantic components stay constant (run as manner), and what changes by construction?
\end{frame}

% ------------------------------------------------------------------------
\begin{frame}{“Construction grammar for semantics” summary}
\small
\begin{itemize}
  \item Meaning is distributed across \textbf{lexemes + constructions + context}.
  \item Constructions range from concrete to abstract; productivity is graded.
  \item Key semantic phenomena:
    \begin{itemize}
      \item \textbf{constructional meaning} (transfer, caused motion, result state)
      \item \textbf{coercion} (non-motion verbs in motion construals)
      \item \textbf{frame alignment} (roles, inferences)
      \item \textbf{discourse meaning} (topic/focus packaging)
    \end{itemize}
  \item A constructicon is a \textbf{network} of related form--meaning pairings.
\end{itemize}
\end{frame}

% ------------------------------------------------------------------------
\begin{frame}{Two quick activities (examples-heavy)}
\small
\textbf{Activity A (classification):} Which construction is at work?

\begin{exe}
  \ex \emph{He talked his way out of trouble.}
  \ex \emph{She sang the baby to sleep.}
  \ex \emph{They emailed the form to the office.}
  \ex \emph{It was Maria who solved the puzzle.}
\end{exe}


\vspace{0.5em}
\textbf{Activity B (productivity test):} Acceptable? What meaning do you get?

\begin{exe}
  \ex \emph{He yawned the cat off the sofa.}
  \ex \emph{She smiled her way into the committee.}
  \ex \emph{They whispered him the secret.}
\end{exe}

\end{frame}

% ------------------------------------------------------------------------
\begin{frame}{Closing: one slide “cheat sheet”}
\small
\begin{tabular}{@{}p{0.32\linewidth}p{0.64\linewidth}@{}}
\toprule
\textbf{Concept} & \textbf{One-liner} \\
\midrule
Construction & Conventional pairing of form + meaning (+ constraints) \\
Constructicon & Inventory/network of constructions in a language \\
Constructional meaning & Meaning contributed by the pattern (not only the verb) \\
Coercion & Construction forces a compatible construal when lexeme alone wouldn’t \\
Inheritance network & Generalizations captured as links among constructions \\
Usage-based learning & Frequency + analogy shape what is stored and what generalizes \\
\bottomrule
\end{tabular}

\vspace{0.75em}
\textbf{Last example to remember:} \emph{whistle your way to work} $\Rightarrow$ motion meaning without a motion verb.
\end{frame}

%\input{part/project-mwe.tex}

\begin{frame}{Conclusions}
  \begin{itemize}
  \item We most often combine words compositionally
  \item We can combine words in non-compositional ways
  \item But still there are shared conventions
  \item Some things we have to learn (\txx{idioms})
  \item The distinction is somewhat fuzzy
  \item A good reference is \citet{Sag02a}
  \item A good reference for Czech Idioms is
    \url{https://www.ceskezvyky.cz/}
\end{itemize}
\end{frame}



\begin{frame}[allowframebreaks]
  \frametitle{References}
  \bibliographystyle{acl_natbib}
  \bibliography{abb,mtg,ling,nlp}
\end{frame}




\end{document}

%%% Local Variables: 
%%% coding: utf-8
%%% mode: latex
%%% TeX-PDF-mode: t
%%% TeX-engine: xetex
%%% End: 

